\documentclass[12pt, titlepage]{article}

\usepackage{amsmath, mathtools}

\usepackage[round]{natbib}
\usepackage{amsfonts}
\usepackage{amssymb}
\usepackage{graphicx}
\usepackage{colortbl}
\usepackage{hyperref}
\usepackage{longtable}
\usepackage{xfrac}
\usepackage{tabularx}
\usepackage{float}
\usepackage{siunitx}
\usepackage{booktabs}
\usepackage{multirow}
\usepackage[section]{placeins}
\usepackage{caption}
\usepackage{fullpage}

\hypersetup{
colorlinks=true,       % false: boxed links; true: colored links
linkcolor=red,          % color of internal links (change box color with linkbordercolor)
citecolor=blue,      % color of links to bibliography
filecolor=magenta,  % color of file links
urlcolor=cyan          % color of external links
}

\usepackage{array}


\input{../../Comments.text}
\input{../../Common.text}

\newcommand{\pdfbase}{https://github.com/richy1337/Verity/blob/main/pdfs}
\newcommand{\ty}[1]{\textsf{#1}}
\newcommand{\key}{\ty{CheckKey}}
\newcommand{\Const}[1]{\texttt{#1}}
\setlength{\emergencystretch}{2em}
\setcounter{tocdepth}{2}

\begin{document}

\title{Module Interface Specification for Verity}

\author{\authname}

\date{November 21, 2026}

\maketitle

\pagenumbering{roman}

\section{Revision History}

\begin{tabularx}{\textwidth}{p{2.3cm}p{3.2cm}X}
\toprule {\bf Date} & {\bf Developer(s)} & {\bf Change}\\
\midrule
2026-11-16 & Richard Han, Calvin Salsali & Prepared Revision $-1$ defining operations and shared data types for the ten modules, including consent, lifecycle, and cleanup contracts.\\
2026-11-18 & Jaden Moore & Clarified initialization rules, typed state/environment variables, literal constant values, and formal lifecycle/metric semantics.\\
2026-11-21 & Calvin Salsali, Raymond Tao, Jaden Moore, Richard Han & Completed individual design reflections (Prompts 1--6) and updated the Generative AI disclosure with confirmed tool versions and verification records.\\
\bottomrule
\end{tabularx}

\clearpage

\tableofcontents

\listoftables

\clearpage

\pagenumbering{arabic}

\section{Introduction}\label{sec:introduction}

This document specifies the ten modules in the
\href{\pdfbase/Design/SoftArchitecture/MG.pdf}{Module Guide (MG)}. It uses the course Module Interface Specification (MIS) template, based on \citet{HoffmanAndStrooper1995}. Each contract lists its dependencies, state, environment, constants, operation signatures and observable behaviour. We collect shared record types here so the same terminology is used on both sides of an interface.

The \href{\pdfbase/SRS/SRS.pdf}{Software Requirements Specification (SRS)} defines the behaviour and limits. The \href{\pdfbase/HazardAnalysis/HazardAnalysis.pdf}{Hazard Analysis} documents the harms addressed by the interfaces. The \href{\pdfbase/VnVPlan/VnVPlan.pdf}{Verification and Validation Plan} documents the planned tests. The design describes a prospective implementation. FR, NFR, SR and INV denote functional requirements, nonfunctional requirements, safety requirements and invariants. The user interface (UI) consumes these application programming interfaces (APIs).

\section{Notation}\label{sec:notation}

An operation $f(x:T)\to U$ takes a value of type $T$ and returns a value of type $U$. Fields are listed in parentheses. A vertical bar between types denotes alternatives. $\operatorname{Optional}(T)$ is either a value of $T$ or absent, and $\operatorname{Seq}(T)$ is an ordered sequence. $\varnothing$ denotes an empty set or map. Time on a host is measured by a monotonic clock. Compare timestamps only on the same clock unless they have been explicitly aligned. A status transition is written $s\to s'$.

A returned reason is safe to show or log and contains no raw selection. All operations that reject an authorization or state check leave state unchanged unless their contract describes a terminal transition. For an unknown or foreign request, the visible reason is the same so it cannot reveal another session's activity. Routine failures are in the result types below.

\renewcommand{\arraystretch}{1.2}
\begin{longtable}{@{}p{0.25\linewidth}p{0.72\linewidth}@{}}
\caption{Shared interface types}\label{TblTypes}\\
\toprule \textbf{Type} & \textbf{Meaning} \\ \midrule
\endfirsthead
\toprule \textbf{Type} & \textbf{Meaning} \\ \midrule
\endhead
\ty{BrowserDocument}, \ty{RenderSurface} & Browser-owned document and extension-owned display surface. \\
\ty{InputEvent}, \ty{Channel} & User input event and authenticated transport channel. \\
\ty{Clock}, \ty{FileSet}, \ty{Directory}, \ty{Storage} & Monotonic time source, versioned files, output directory and persistent storage. \\
\ty{Session}, \ty{RequestId} & Opaque authenticated session context and a fresh UUIDv4 request identifier (Universally Unique Identifier, version 4). A request identifier alone grants no access. \\
\key & Pair of initiating session identity and request identifier. Services check the session against trusted transport context. \\
\ty{Boolean} & True or false. $\mathbb{N}_0$ denotes nonnegative integers. \\
\ty{ModelArtifact} & Loaded immutable model parameters, tokenizer or parser resources. \\
\ty{Unit} & Successful operation with no data returned. \\
\ty{Modality} & Text or image. \\
\ty{Approve}, \ty{Decline} & Alternatives of an explicit user decision. \\
\ty{WorkEvent} & Tagged worker result, timeout, dependency error or accepted cancellation, with check key and a trusted work capability. Client-supplied events never receive that capability. \\
\ty{RequestRecord} & Status, authenticated owner, admission instant, deadline, approved route and opaque work handle; latest safe \ty{Update}; and optional terminal instant and delivery-expiry instant. \\
\ty{WorkerLease} & Opaque authorization for one worker and check. \\
\ty{LocalInstant} & Timestamp on the cleanup owner's monotonic clock. \\
\ty{CleanupReceipt} & Key, local completion time and cleanup outcome. \\
\ty{MetricReport} & Per-class and per-slice counts and metrics, with eligibility and coverage denominators. \\
\ty{Gesture}, \ty{View} & Explicit browser invocation for one selection, and a trusted extension view scoped to its tab and session. \\
\ty{RawItem} & One text or static image selection with modality, declared media attributes and optional user-entered context. Modality is text or image. \\
\ty{Item} & Bounded text or validated static image with safe descriptors; image content includes the bounded original container for metadata inspection and a safely decoded pixel view. The transient content is accessible only to work authorized for the check. \\
\ty{Route} & Team service, or a configured third-party provider with recipient, transfer scope and current retention terms. \\
\ty{ConsentPreview} & Modality, recipient identity, transfer scope, retention terms, policy version and route eligibility. Contains no raw content. \\
\ty{Consent} & One-time authorization binding session, check key, selected item, intended recipient, scope and policy version. \\
\ty{Job} & Check key, one request-scoped input and selected route with optional consent. Browser UI state and admitted service state are recorded separately. \\
\ty{Reason} & One of invalid input, input unavailable, unsupported host, busy, timeout, dependency failure, unauthorized, weak evidence, view unavailable, invalid state, not initialized, configuration mismatch, cleanup violation, invalid manifest or missing outcomes; with safe explanatory text. Unknown and foreign keys both return unauthorized. \\
\ty{Status} & Pending, completed, inconclusive, failed or cancelled. \\
\ty{Evidence} & Claim, nonempty safe locator and source kind (Observed Input Signal, External Reference or Model Inference). Only supported claims are displayed. \\
\ty{Signals} & Candidate evidence, optional score in $[0,1]$, model version, modality, required-work availability, applicability (supported, unsupported or uncertain), and optional separate \ty{Provenance}. \\
\ty{Provenance} & Separate EXIF presence/structure and C2PA verification observations. Each records present and valid, absent or stripped, or inconsistent, plus the check performed. Absent or stripped means absence observed, not proof of removal. EXIF presence is not cryptographic authenticity. Unsupported verification is unavailable, not valid; a parser outage is a dependency condition. \\
\ty{Assessment} & Optional direction (higher or lower synthetic/manipulation signal), evidence, visible uncertainty, limits and a next check. \\
\ty{Update} & Check key, status, and optional assessment or reason. The status determines the variant: pending has no assessment; completed requires an assessment with direction and supported evidence; inconclusive requires limits and guidance; failed/cancelled require a reason and have no direction. Inconclusive updates can have supported observations and have no direction. \\
\ty{Decision} & Completed with an assessment, or inconclusive with limitations and guidance. Includes no job or session state. \\
\ty{Config} & Module-specific immutable configuration. Host settings include trusted endpoint and supported origins; coordinator settings include capacity and worker configuration; buffer settings include cleanup policy; offline settings include output directory, seed and model profiles. Detector/decision settings contain versioned task and modality; frozen model, applicability-gate and calibration profile identifiers; score direction; abstention margin; supported domain; and required signal set. Runtime and offline settings are separate instances. An absent genre profile cannot assert supported applicability. \\
\ty{LeaseRecord} & Authorized worker capability, transient item references and open/closed marker for one check. \\
\ty{ViewRecord} & Active key, local status, cancellation marker and safe display model. \\
\ty{Manifest}, \ty{Run} & Labelled corpus revision, item identity and split metadata. Run adds frozen settings, outcomes, audit and metric report. \\
\bottomrule
\end{longtable}

\noindent The numeric limits below are exported from the relevant modules. SRS values are reproduced here; the delivery-record bound is a design decision. In the contracts, a caller depends on the named limit. The service processing deadline, UI display deadline and cleanup interval measure three distinct events. The delivery-record age is a design bound for safe status metadata only: one minute permits transient transport retries while bounding storage; it never permits raw content retention.

\begin{longtable}{@{}p{0.56\linewidth}p{0.38\linewidth}@{}}
\caption{Exported constant values (SRS unless marked design)}\label{TblConstants}\\
\toprule Constant & Value and unit \\ \midrule
\endfirsthead
\toprule Constant & Value and unit \\ \midrule
\endhead
\Const{MIN\_\allowbreak{}TEXT\_\allowbreak{}LENGTH} & 150 Unicode code points \\
\Const{MAX\_\allowbreak{}TEXT\_\allowbreak{}LENGTH} & 10,000 Unicode code points \\
\Const{MAX\_\allowbreak{}IMAGE\_\allowbreak{}SIZE} & 10.0 MB ($10{,}000{,}000$ bytes) \\
\Const{MIN\_\allowbreak{}IMAGE\_\allowbreak{}DIM} & $64\times64$ pixels, inclusive per axis \\
\Const{MAX\_\allowbreak{}IMAGE\_\allowbreak{}DIM} & $4096\times4096$ pixels, inclusive per axis \\
\Const{MAX\_\allowbreak{}RESPONSE\_\allowbreak{}TIME} & 10.0 seconds from service admission \\
\Const{MAX\_\allowbreak{}CONCURRENT\_\allowbreak{}REQS} & 10 isolated admitted requests \\
\Const{MAX\_\allowbreak{}DELIVERY\_\allowbreak{}RECORD\_\allowbreak{}AGE} & 60 seconds after terminal decision (design) \\
\Const{MAX\_\allowbreak{}RETENTION\_\allowbreak{}TIME} & 1.0 second after local terminal decision \\
\Const{MAX\_\allowbreak{}UI\_\allowbreak{}RENDER\_\allowbreak{}LATENCY} & 500 milliseconds \\
\Const{MAX\_\allowbreak{}CANCEL\_\allowbreak{}UI\_\allowbreak{}LATENCY} & 200 milliseconds \\
\Const{MIN\_\allowbreak{}F1\_\allowbreak{}SCORE} & 0.80 assessed macro-average $F_1$ \\
\Const{MIN\_\allowbreak{}COVERAGE} & 0.70 of all eligible items \\
\Const{MAX\_\allowbreak{}FPR} & 0.10 for assessed and full-negative accounting \\
\Const{MIN\_\allowbreak{}BASELINE\_\allowbreak{}DELTA} & $+0.10$ assessed macro-average $F_1$ \\
\Const{MAX\_\allowbreak{}REPRO\_\allowbreak{}DELTA} & 0.005 absolute macro-average $F_1$ difference \\
\bottomrule
\end{longtable}

\textbf{Wire-contract preservation.} The shared records abstract the SRS S.5 transport envelope; they do not replace it. M1 preserves the required request fields \texttt{req\_id}, \texttt{session\_token}, \texttt{modality}, \texttt{payload} and \texttt{consent\_granted}. It preserves the response fields \texttt{req\_id}, \texttt{status}, \texttt{signals[]}, \texttt{uncertainty}, \texttt{limitations[]} and \texttt{verification\_\allowbreak{}steps[]}. Optional user context remains part of the explicitly selected transfer scope. Evidence source kinds serialize as \texttt{observed\_\allowbreak{}signal}, \texttt{external\_\allowbreak{}reference} or \texttt{inference}; their descriptive labels are display text. M5 checks session ownership against trusted transport context rather than trusting a payload field.

\textbf{Initialization and failures.} M4 is a stateless library. M1--M3 and M5--M10 are abstract objects: one registry per owning extension/service process, and one M10 object per offline run. In particular, M2 and M3 keep maps across views rather than creating an implicit global active tab. Each abstract object exports the common operation
\[
\texttt{initialize}(p:\ty{Config})\to\ty{Unit}\mid\ty{Reason}.
\]
Its Boolean state $\mathit{ready}$ is initially false. Successful initialization establishes empty maps and immutable configuration $p$, then sets $\mathit{ready}'=\mathrm{true}$ and $out:=\ty{Unit}$. Invalid or unavailable configuration gives $out:=\ty{Reason}$ and leaves $\mathit{ready}=\mathrm{false}$. Reinitialization returns invalid state without replacing active work. Except M4, every operation requires $\mathit{ready}$; otherwise $out:=\text{not initialized}$ with no state change. Model loading occurs here, not when handling a request. For M1--M3 the configuration additionally contains the trusted endpoint/host settings; M10 configuration contains its output directory, fixed at initialization. Configuration contains no user payload. Expected failures use the declared Reason alternative; $exc:=\text{none}$ in every contract. State-changing operations return a receipt or acknowledgment of their effect where documented below. All preserve ownership and terminal-state absorption.

\textbf{State convention.} $X'$ denotes post-state and $out$ the returned value. Each abstract object also has the shared Boolean ready and immutable configuration state from initialization. Loaded model parameters and calibration tables are in-memory state of M7, M8 and M9 respectively; configuration identifiers name those artifacts rather than substituting for their loaded values. For operations receiving a configuration argument, it must match the initialized profile; mismatch returns configuration mismatch without model reload. An unmentioned state or environment variable is unchanged. Empty initial maps satisfy their ownership invariants. For every successful operation, the output has the declared type and the stated invariants hold; rejected authorization never discloses another key's state. The operation descriptions below define the output/effect beyond these common rules.

\textbf{Data ownership.} Each field has one module owner. M1 owns host authorization metadata; M2 owns the unsent selected item and consent until dispatch; M3 owns safe display state; M5 owns service lifecycle and delivery metadata; M6 owns admitted raw-content buffers and worker references; M7--M9 own only loaded immutable analysis artifacts between calls; M10 owns the offline run. A module may read another's value through its interface without owning a second mutable copy. Call-local input/output values are temporary borrows, not additional module state. Dispatch transfers staged payload ownership from M2 to the service's M6 and clears the browser staging reference; transport copies are request-owned work released on completion, cancellation or error. M1 and M4 retain no content after their calls. The selected webpage remains external browser state, outside application-owned buffers. M6 cleanup controls its own process's buffers; M2 releases browser staging independently when its local check ends. Safe display records may remain visible and contain no raw or reconstructive excerpts.

\textbf{Operation effects.} Read/query operations return data without changing module state; command operations return only Unit on success except the documented acknowledgments. M1 capture and transport, M2 begin/authorize, M5 admission/advance/cancel, M6 hold/release and M10 evaluate/recompute return a key, receipt, safe update or recorded run as acknowledgment of an external or state-changing action. These exceptions avoid a second query racing with that action. Expected failures always return Reason; no exception carries a raw item.

\section{Module Decomposition}

The hierarchy assigns one specification to each leaf module in the MG.

\begin{table}[ht]
\centering
\begin{tabular}{p{0.34\textwidth} p{0.56\textwidth}}
\toprule
\textbf{Level 1} & \textbf{Level 2}\\
\midrule

Hardware-Hiding & M1: Browser and host bridge \\
\midrule

Behaviour-Hiding & M2: Selection and consent\\
& M3: Result card\\
& M4: Input validation\\
& M5: Request coordinator\\
& M6: Temporary content\\
\midrule

Software-Decision-Hiding & M7: Text signals\\
& M8: Image signals and provenance\\
& M9: Assessment policy\\
& M10: Offline evaluation\\
\bottomrule

\end{tabular}
\caption{Module Hierarchy}
\label{TblMH}
\end{table}

\clearpage

%% ------------------------------------------------------------------ M1
\section{MIS of M1 Browser and host bridge} \label{sec:m1}

M1 isolates browser access and the authenticated service connection from its callers.

\subsection{Module}

Browser and host bridge

\subsection{Uses}

M4 (exported input limits).

\subsection{Syntax}

The following constants and operations form the public interface. Shared initialization is defined in Section~\ref{sec:notation}.

\subsubsection{Exported Constants}

None.

\subsubsection{Exported Access Programs}

\begin{center}
\begin{tabular}{p{3.3cm} p{4cm} p{4cm} p{2.4cm}}
\hline
\textbf{Name} & \textbf{In} & \textbf{Out} & \textbf{Exceptions} \\
\hline
\texttt{initialize} & $p:\ty{Config}$ & $\ty{Unit}\mid\ty{Reason}$ & None \\
\texttt{readSelection} & $g:\ty{Gesture}$ & $\ty{RawItem}\mid\ty{Reason}$ & None \\
\texttt{send} & $s:\ty{Session},\ j:\ty{Job}$ & $\ty{Update}\mid\ty{Reason}$ & None \\
\texttt{sendCancel} & $s:\ty{Session},\ k:\key$ & $\ty{Update}\mid\ty{Reason}$ & None \\
\texttt{draw} & $v:\ty{View},\ u:\ty{Update}\mid\ty{Reason}$ & $\ty{Unit}\mid\ty{Reason}$ & None \\
\hline
\end{tabular}
\end{center}

\subsection{Semantics}

The state and environment below are private; callers rely on the operation effects.

\subsubsection{State Variables}

$G:\ty{View}\rightharpoonup\ty{Gesture}$ records current capture authorization and is initially empty. No gesture authorizes a different tab or document.

\subsubsection{Environment Variables}

\begin{itemize}
\item $\mathit{page}:\ty{BrowserDocument}$: the webpage open in the user's tab, including the user's current text or image selection. Read by \texttt{readSelection}; never modified. Navigating away invalidates the current gesture.
\item $\mathit{gesture}:\ty{Gesture}$: the user's explicit browser invocation (context menu or keyboard shortcut) on the selection. Read by \texttt{readSelection}.
\item $\mathit{screen}:\ty{RenderSurface}$: the extension-controlled view region of the browser window. Changed by \texttt{draw}, which replaces the displayed content with the rendered update or reason.
\item $\mathit{link}:\ty{Channel}$: the encrypted connection to the configured team service. Changed by \texttt{send} and \texttt{sendCancel}, which transmit one command and receive its reply.
\item $\mathit{clock}:\ty{Clock}$: the browser host's monotonic clock. Read only.
\end{itemize}

\subsubsection{Access Routine Semantics}

\noindent \texttt{readSelection}():
Check that the gesture is current, user initiated and scoped to its tab. Read only its chosen item. Preliminary capture-size checks use M4's exported input limits; they never replace M4 validation at service admission. If the browser cannot provide the selected item within the authorized scope and input bounds, return input unavailable. An expired or restricted gesture returns a corrective reason. A change of page invalidates the gesture authorization.

\medskip\noindent \texttt{send}(), \texttt{sendCancel}():
Carry the command over the authenticated encrypted connection. A response belongs to the same initiating context. Return a transport reason on connection failure and never retry a content transfer using a different route.

\medskip\noindent \texttt{draw}():
Render within the extension-controlled view. Treat all text as inert. If that view has been closed or its context has changed, return a view-unavailable reason. The host does not decide the active check or supply a verdict.

%% ------------------------------------------------------------------ M2
\clearpage
\section{MIS of M2 Selection and consent} \label{sec:m2}

M2 binds one scoped selection and its consent to a fresh check.

\subsection{Module}

Selection and consent

\subsection{Uses}

M1.

\subsection{Syntax}

The following constants and operations form the public interface. Shared initialization is defined in Section~\ref{sec:notation}.

\subsubsection{Exported Constants}

None.

\subsubsection{Exported Access Programs}

\begin{center}
\begin{tabular}{p{3.3cm} p{4cm} p{4cm} p{2.4cm}}
\hline
\textbf{Name} & \textbf{In} & \textbf{Out} & \textbf{Exceptions} \\
\hline
\texttt{initialize} & $p:\ty{Config}$ & $\ty{Unit}\mid\ty{Reason}$ & None \\
\texttt{begin} & $g:\ty{Gesture}$ & $\key\mid\ty{Reason}$ & None \\
\texttt{preview} & $k:\key$ & $\ty{ConsentPreview}\mid\ty{Reason}$ & None \\
\texttt{authorize} & $k:\key,\ a:\ty{Approve}\mid\ty{Decline}$ & $\ty{Update}\mid\ty{Reason}$ & None \\
\texttt{discard} & $k:\key$ & $\ty{Unit}\mid\ty{Reason}$ & None \\
\hline
\end{tabular}
\end{center}

\subsection{Semantics}

The state and environment below are private; callers rely on the operation effects.

\subsubsection{State Variables}

$Q:\ty{View}\rightharpoonup\ty{Job}$ stores at most one staged job and optional consent per view. Initially $Q=\varnothing$. A consent in $Q$ binds exactly its job key, item and route.

\subsubsection{Environment Variables}

\begin{itemize}
\item $\mathit{consentPrompt}:\ty{RenderSurface}$: the consent disclosure shown on screen for a proposed third-party route. Changed by \texttt{preview}, which displays recipient, scope and retention terms, and cleared by \texttt{authorize} and \texttt{discard}.
\item $\mathit{userDecision}:\ty{InputEvent}$: the user's Approve or Decline action entered with keyboard or mouse. Read by \texttt{authorize}.
\end{itemize}

\subsubsection{Access Routine Semantics}

\noindent \texttt{begin}():
Read the item through M1, replace any unsent selection in the same view and obtain a fresh UUIDv4 check key bound to its authenticated session. Make its modality and bounds available before processing. Reject restricted, missing or out-of-scope selections without a provider call.

\medskip\noindent \texttt{preview}():
Return a ConsentPreview for that key, including its modality and route eligibility. When a third party is proposed, include recipient, content scope, retention rules and the limit on recalling already delivered data. Unknown or incompatible terms disable that route.

\medskip\noindent \texttt{authorize}():
A team-service check proceeds on the user's explicit check action. A third-party check requires explicit approval after its preview. Decline clears the unsent item and returns a local cancelled update without admitting service work or sending content. Approval for the active key authorizes exactly the disclosed item, scope and recipient. Dispatch once through M1, transfer payload ownership and clear the staging reference as described in Section~\ref{sec:notation}. An expired decision cannot authorize a later gesture. Duplicate authorization after dispatch returns invalid state and sends nothing; item, route or policy-version changes require a new preview and approval.

\medskip\noindent \texttt{discard}():
Clear an unsent selection and its authorization. An admitted check is cancelled through the result card and request coordinator.

%% ------------------------------------------------------------------ M3
\clearpage
\section{MIS of M3 Result card} \label{sec:m3}

M3 owns the active result binding and uncertainty presentation for each view.

\subsection{Module}

Result card

\subsection{Uses}

M1.

\subsection{Syntax}

The following constants and operations form the public interface. Shared initialization is defined in Section~\ref{sec:notation}.

\subsubsection{Exported Constants}

\Const{MAX\_\allowbreak{}UI\_\allowbreak{}RENDER\_\allowbreak{}LATENCY}, \Const{MAX\_\allowbreak{}CANCEL\_\allowbreak{}UI\_\allowbreak{}LATENCY}.

\subsubsection{Exported Access Programs}

\begin{center}
\begin{tabular}{p{3.3cm} p{4cm} p{4cm} p{2.4cm}}
\hline
\textbf{Name} & \textbf{In} & \textbf{Out} & \textbf{Exceptions} \\
\hline
\texttt{initialize} & $p:\ty{Config}$ & $\ty{Unit}\mid\ty{Reason}$ & None \\
\texttt{start} & $v:\ty{View},\ k:\key$ & $\ty{Unit}\mid\ty{Reason}$ & None \\
\texttt{accept} & $v:\ty{View},\ u:\ty{Update}$ & $\ty{Unit}\mid\ty{Reason}$ & None \\
\texttt{reject} & $v:\ty{View},\ k:\key,\ r:\ty{Reason}$ & $\ty{Unit}\mid\ty{Reason}$ & None \\
\texttt{cancel} & $v:\ty{View},\ k:\key$ & $\ty{Unit}\mid\ty{Reason}$ & None \\
\hline
\end{tabular}
\end{center}

\subsection{Semantics}

The state and environment below are private; callers rely on the operation effects.

\subsubsection{State Variables}

$V:\ty{View}\rightharpoonup\ty{ViewRecord}$, initially empty. $V(v)$ contains only the active key for view $v$. Once locally cancelled, that key cannot receive an assessment; changing the active key clears all prior evidence.

\subsubsection{Environment Variables}

\begin{itemize}
\item $\mathit{card}:\ty{RenderSurface}$: the result card region on screen. Changed by \texttt{start} (clears previous evidence, shows pending after admission), \texttt{accept} (shows the update), \texttt{reject} (shows the corrective reason) and \texttt{cancel} (shows cancelled).
\item $\mathit{keyboard}:\ty{InputEvent}$: keyboard and pointer input on the card's controls, such as the cancel control. Read only.
\item $\mathit{clock}:\ty{Clock}$: the browser host's monotonic clock, used to meet \Const{MAX\_\allowbreak{}UI\_\allowbreak{}RENDER\_\allowbreak{}LATENCY} and \Const{MAX\_\allowbreak{}CANCEL\_\allowbreak{}UI\_\allowbreak{}LATENCY}. Read only.
\end{itemize}

\subsubsection{Access Routine Semantics}

\noindent \texttt{start}():
Replace the previous view binding with the new check and clear the previous evidence. Show pending only after admission is acknowledged. A terminal reply that reaches the active view before the pending update may render directly.

\medskip\noindent \texttt{accept}():
Require the full active key, an open view and an allowed transition. Ignore mismatched, stale and post-terminal replies. Display a valid update by \Const{MAX\_\allowbreak{}UI\_\allowbreak{}RENDER\_\allowbreak{}LATENCY} of receipt. Show source kind and locator for evidence, visible uncertainty, plain-language limits and a next check. Use a text label and keyboard-accessible controls. A higher-signal result warns that an authentic item can be flagged. A lower-risk label requires a valid completed update with the lower-signal direction. A malformed completed result is displayed as a local failed check, with no verdict; it does not rewrite an already terminal service record. An active pre-admission rejection renders by \Const{MAX\_\allowbreak{}UI\_\allowbreak{}RENDER\_\allowbreak{}LATENCY} from rejection receipt.

\medskip\noindent \texttt{reject}():
Display a corrective reason for an active pre-admission rejection. There is no service status or detector result. $out:=\ty{Unit}$ after rendering; unmatched keys leave the view unchanged.

\medskip\noindent \texttt{cancel}():
For an active pending check, immediately set the local cancellation marker, remove assessment fields, show cancelled by \Const{MAX\_\allowbreak{}CANCEL\_\allowbreak{}UI\_\allowbreak{}LATENCY}, and request service cancellation through M1. Service receipt and cleanup confirmation are separate. Suppress all later assessments for that view and key even if a result was already on its way.

%% ------------------------------------------------------------------ M4
\clearpage
\section{MIS of M4 Input validation} \label{sec:m4}

M4 admits only bounded inert text or supported static image content.

\subsection{Module}

Input validation

\subsection{Uses}

No design module.

\subsection{Syntax}

The following constants and operations form the public interface. Shared initialization is defined in Section~\ref{sec:notation}.

\subsubsection{Exported Constants}

\Const{MIN\_\allowbreak{}TEXT\_\allowbreak{}LENGTH}, \Const{MAX\_\allowbreak{}TEXT\_\allowbreak{}LENGTH}, \Const{MAX\_\allowbreak{}IMAGE\_\allowbreak{}SIZE},\newline \Const{MIN\_\allowbreak{}IMAGE\_\allowbreak{}DIM}, \Const{MAX\_\allowbreak{}IMAGE\_\allowbreak{}DIM}.

\subsubsection{Exported Access Programs}

\begin{center}
\begin{tabular}{p{3.3cm} p{4cm} p{4cm} p{2.4cm}}
\hline
\textbf{Name} & \textbf{In} & \textbf{Out} & \textbf{Exceptions} \\
\hline
\texttt{validate} & $x:\ty{RawItem}$ & $\ty{Item}\mid\ty{Reason}$ & None \\
\hline
\end{tabular}
\end{center}

\subsection{Semantics}

The state and environment below are private; callers rely on the operation effects.

\subsubsection{State Variables}

None.

\subsubsection{Environment Variables}

None. The module has no interaction outside the software.

\subsubsection{Access Routine Semantics}

\noindent \texttt{validate}():
For text, require valid UTF-8 and count Unicode code points, not UTF-8 bytes or JavaScript UTF-16 code units. Apply inclusive minimum and maximum bounds. Keep the selection as inert text. For an image, use decimal MB ($10^6$ bytes) and check the byte bound before decoding, require a supported static Portable Network Graphics (PNG), Joint Photographic Experts Group (JPEG) or WebP format, compare headers with the decoded type, and bound both dimensions and decoding work. Allow only the selected item and the explicitly supplied context. Reject empty, out-of-bounds, inconsistent, unsupported or corrupt content with a specific reason and release temporary decode references. No detector is invoked on rejection. $out:=\ty{Item}$ on success, otherwise $out:=\ty{Reason}$; no module state changes.

%% ------------------------------------------------------------------ M5
\clearpage
\section{MIS of M5 Request coordinator} \label{sec:m5}

M5 serializes request admission and terminal decisions while preserving session isolation.

\subsection{Module}

Request coordinator

\subsection{Uses}

M4, M6--M9.

\subsection{Syntax}

The following constants and operations form the public interface. Shared initialization is defined in Section~\ref{sec:notation}.

\subsubsection{Exported Constants}

\Const{MAX\_\allowbreak{}RESPONSE\_\allowbreak{}TIME}, \Const{MAX\_\allowbreak{}CONCURRENT\_\allowbreak{}REQS}, \Const{MAX\_\allowbreak{}DELIVERY\_\allowbreak{}RECORD\_\allowbreak{}AGE}.

\subsubsection{Exported Access Programs}

\begin{center}
\begin{tabular}{p{3.3cm} p{4cm} p{4cm} p{2.4cm}}
\hline
\textbf{Name} & \textbf{In} & \textbf{Out} & \textbf{Exceptions} \\
\hline
\texttt{initialize} & $p:\ty{Config}$ & $\ty{Unit}\mid\ty{Reason}$ & None \\
\texttt{admit} & $s:\ty{Session},\ j:\ty{Job}$ & $\ty{Update}\mid\ty{Reason}$ & None \\
\texttt{advance} & $k:\key,\ e:\ty{WorkEvent}$ & $\ty{Update}\mid\ty{Reason}$ & None \\
\texttt{cancel} & $s:\ty{Session},\ k:\key$ & $\ty{Update}\mid\ty{Reason}$ & None \\
\texttt{status} & $s:\ty{Session},\ k:\key$ & $\ty{Update}\mid\ty{Reason}$ & None \\
\hline
\end{tabular}
\end{center}

\subsection{Semantics}

The state and environment below are private; callers rely on the operation effects.

\subsubsection{State Variables}

Let $R:\key\rightharpoonup\ty{RequestRecord}$ be a partial map. Each record contains the RequestRecord fields from Table~\ref{TblTypes}, including its latest safe update and optional terminal/expiry times. The update is pending at admission and is replaced atomically with each status decision. Only the terminal/expiry times are absent until the first terminal decision; the expiry instant is then fixed to terminal time plus the delivery-record age. Initially $R=\varnothing$. For every $k\in\operatorname{dom}(R)$, the record owner equals the session component of $k$. Let $P$ count pending records and $K$ be the configured admission capacity; $0\le P\le K$ and $K\ge\Const{MAX\_\allowbreak{}CONCURRENT\_\allowbreak{}REQS}$. A terminal record never changes status.

\subsubsection{Environment Variables}

\begin{itemize}
\item $\mathit{clock}:\ty{Clock}$: the service host's monotonic clock. Read by \texttt{admit} to set the deadline and by \texttt{advance} to decide whether the deadline is reached. Never modified.
\end{itemize}

\subsubsection{Access Routine Semantics}

\noindent \texttt{admit}():
Check the owner, key freshness, current consent and capacity, then validate through M4. Discard rejected work before entry into the state map. Consume each provider consent once and bind it to the validated item and provider policy. For success, register transient ownership in M6, enter pending, set the deadline to admission time plus \Const{MAX\_\allowbreak{}RESPONSE\_\allowbreak{}TIME}, and dispatch to the correct detector. Deploy enough capacity for the SRS concurrency target. Refuse excess pre-admission work with a busy reason. Never substitute a provider for the approved route.

\medskip\noindent \texttt{advance}():
Only the check's trusted worker, cancellation handler or timer can deliver a work event. While pending, let $c$ mean accepted cancellation, $d$ mean deadline reached, $e$ mean unrecoverable dependency error, and $a$ mean required work has finished or been declared unavailable. Let $v$ be a total Boolean that is false until $a$ and true exactly when M9 supplies a supported completed assessment. An available result from M9 is either completed or inconclusive. Evaluate the following in order as one atomic decision.

\begin{table}[H]
\centering
\caption{Atomic lifecycle decision, matching SRS S.9.1}\label{TblLifecycle}
\begin{tabularx}{0.92\linewidth}{@{}lX@{}}
\toprule Condition & Next state \\ \midrule
$c$ & cancelled \\
$\neg c \land (d\lor e)$ & failed, with the appropriate reason \\
$\neg c \land \neg d \land \neg e \land \neg a$ & pending \\
$\neg c \land \neg d \land \neg e \land a \land \neg v$ & inconclusive \\
$\neg c \land \neg d \land \neg e \land a \land v$ & completed \\
\bottomrule
\end{tabularx}
\end{table}

The five rows are disjoint and total over all 32 assignments of $(c,d,e,a,v)$; UT-C4 in the V\&V Plan specifies exhaustive verification of them. Partial unavailable work without an unrecoverable error sets $a$ when resolved and $v=\mathrm{false}$, producing inconclusive. Once terminal, the status is fixed. On the first terminal decision, stop request-owned work and invoke M6 cleanup. Retain only the safe response needed for delivery. Unrecognized work events and late updates cannot mutate the state. A foreign key exposes no state.

\medskip\noindent \texttt{cancel}():
Check the authenticated owner. For pending work, submit an accepted cancellation event to the same decision as \texttt{advance}. A terminal request remains terminal. Duplicate cancellation is safe.

\medskip\noindent \texttt{status}():
Return the stored safe update only to the requesting session. A timer retires each terminal record at its stored expiry instant, \Const{MAX\_\allowbreak{}DELIVERY\_\allowbreak{}RECORD\_\allowbreak{}AGE} after the terminal decision; subsequent queries return the same generic reason as for an unknown key. The record contains no raw item or reconstructive evidence excerpt. Expiry is scheduled at the terminal decision and does not depend on a later status poll. The status query itself leaves state unchanged; no delivery-acknowledgment operation is required.

%% ------------------------------------------------------------------ M6
\clearpage
\section{MIS of M6 Temporary content} \label{sec:m6}

M6 controls worker access to transient content and records its logical release.

\subsection{Module}

Temporary content

\subsection{Uses}

No design module.

\subsection{Syntax}

The following constants and operations form the public interface. Shared initialization is defined in Section~\ref{sec:notation}.

\subsubsection{Exported Constants}

\Const{MAX\_\allowbreak{}RETENTION\_\allowbreak{}TIME}.

\subsubsection{Exported Access Programs}

\begin{center}
\begin{tabular}{p{3.3cm} p{4cm} p{4cm} p{2.4cm}}
\hline
\textbf{Name} & \textbf{In} & \textbf{Out} & \textbf{Exceptions} \\
\hline
\texttt{initialize} & $p:\ty{Config}$ & $\ty{Unit}\mid\ty{Reason}$ & None \\
\texttt{hold} & $k:\key,\ x:\ty{Item}$ & $\ty{WorkerLease}\mid\ty{Reason}$ & None \\
\texttt{borrow} & $k:\key,\ w:\ty{WorkerLease}$ & $\ty{Item}\mid\ty{Reason}$ & None \\
\texttt{release} & $k:\key,\ t:\ty{LocalInstant}$ & $\ty{CleanupReceipt}\mid\ty{Reason}$ & None \\
\hline
\end{tabular}
\end{center}

\subsection{Semantics}

The state and environment below are private; callers rely on the operation effects.

\subsubsection{State Variables}

$B:\key\rightharpoonup\ty{LeaseRecord}$, initially empty. A live lease reads only its own key. Closed records contain no item references; late workers cannot recreate a closed key. Closed-key rejection persists through the safe delivery-record lifetime without retaining content.

\subsubsection{Environment Variables}

\begin{itemize}
\item $\mathit{disk}:\ty{Storage}$: persistent storage and log files on the browser host and service host. No access program changes $\mathit{disk}$ with raw payload content; it is unchanged by \texttt{hold}, \texttt{borrow} and \texttt{release}.
\item $\mathit{clock}:\ty{Clock}$: the cleanup owner's monotonic clock. Read by \texttt{release} to check the cleanup bound \Const{MAX\_\allowbreak{}RETENTION\_\allowbreak{}TIME}.
\end{itemize}

\subsubsection{Access Routine Semantics}

\noindent \texttt{hold}():
Register one validated item for a fresh key and return its unforgeable worker lease. M5 delegates that lease only to request-owned work and calls borrow before passing the item to a detector. Formally $B'(k)=\operatorname{open}(x,w)$ and $out:=w$; a closed or occupied key returns invalid state. This registry, workers and exception handlers never persist raw payloads.

\medskip\noindent \texttt{borrow}():
Require a live lease authorized only for the same key. Expose data only for the lease duration and track references owned by that work. A missing, foreign or revoked lease returns the same reason.

\medskip\noindent \texttt{release}():
Revoke leases, stop request-owned work and drop application-owned service references, including queued copies, by \Const{MAX\_\allowbreak{}RETENTION\_\allowbreak{}TIME} after the local terminal decision. The argument $t$ is the recorded local terminal-decision instant; measure the bound from $t$, never from the later cleanup-call time. Set $B'(k)$ to a closed record with no item references. Return a receipt only after cleanup. Repeating release leaves the key closed. The receipt covers application-owned references and worker exit. Report a cleanup violation if the bound is missed. The operation cannot recall content that has already reached an approved external recipient.

%% ------------------------------------------------------------------ M7
\clearpage
\section{MIS of M7 Text signals} \label{sec:m7}

M7 produces versioned text signals with explicit domain applicability.

\subsection{Module}

Text signals

\subsection{Uses}

No design module.

\subsection{Syntax}

The following constants and operations form the public interface. Shared initialization is defined in Section~\ref{sec:notation}.

\subsubsection{Exported Constants}

None.

\subsubsection{Exported Access Programs}

\begin{center}
\begin{tabular}{p{3.3cm} p{4cm} p{4cm} p{2.4cm}}
\hline
\textbf{Name} & \textbf{In} & \textbf{Out} & \textbf{Exceptions} \\
\hline
\texttt{initialize} & $p:\ty{Config}$ & $\ty{Unit}\mid\ty{Reason}$ & None \\
\texttt{extractText} & $x:\ty{Item},\ p:\ty{Config}$ & $\ty{Signals}\mid\ty{Reason}$ & None \\
\hline
\end{tabular}
\end{center}

\subsection{Semantics}

The state and environment below are private; callers rely on the operation effects.

\subsubsection{State Variables}

$p_7:\ty{Config}$ is immutable after initialization; it identifies the loaded text detector, attribution strategy and applicability gate. $D_7:\operatorname{Optional}(\ty{ModelArtifact})$ holds loaded detector and applicability resources, absent before initialization. No request content is retained between calls.

\subsubsection{Environment Variables}

\begin{itemize}
\item $\mathit{modelFiles}:\ty{FileSet}$: the approved, versioned text model files on the service host's disk. Read when the detector is loaded; never modified.
\end{itemize}

\subsubsection{Access Routine Semantics}

\noindent \texttt{extractText}():
Require an admitted text item and a compatible frozen detector. Return versioned statistical observations, model inferences and source locators. Treat item and optional context as data throughout. For a known unsupported domain, report applicability as unsupported; for an undecidable domain, report uncertain and withhold a directional score when it cannot be justified. For model or runtime failure, return a dependency reason. The result carries no raw submission or request state. Applicability is supported only if the frozen, version-matched genre gate supports the selected domain; a declared opinion/satire context, unavailable genre profile or uncertain gate result prevents directional assessment. Enable directional text assessment only after the versioned gate passes the held-out genre verification in the V\&V Plan.

%% ------------------------------------------------------------------ M8
\clearpage
\section{MIS of M8 Image signals and provenance} \label{sec:m8}

M8 produces image signals with separately interpreted provenance observations.

\subsection{Module}

Image signals and provenance

\subsection{Uses}

No design module.

\subsection{Syntax}

The following constants and operations form the public interface. Shared initialization is defined in Section~\ref{sec:notation}.

\subsubsection{Exported Constants}

None.

\subsubsection{Exported Access Programs}

\begin{center}
\begin{tabular}{p{3.3cm} p{4cm} p{4cm} p{2.4cm}}
\hline
\textbf{Name} & \textbf{In} & \textbf{Out} & \textbf{Exceptions} \\
\hline
\texttt{initialize} & $p:\ty{Config}$ & $\ty{Unit}\mid\ty{Reason}$ & None \\
\texttt{extractImage} & $x:\ty{Item},\ p:\ty{Config}$ & $\ty{Signals}\mid\ty{Reason}$ & None \\
\hline
\end{tabular}
\end{center}

\subsection{Semantics}

The state and environment below are private; callers rely on the operation effects.

\subsubsection{State Variables}

$p_8:\ty{Config}$ is immutable after initialization and identifies the visual model and provenance parsers. $D_8:\operatorname{Optional}(\ty{ModelArtifact})$ holds loaded resources, absent before initialization. No request content is retained between calls.

\subsubsection{Environment Variables}

\begin{itemize}
\item $\mathit{modelFiles}:\ty{FileSet}$: the approved, versioned image model and provenance parser files on the service host's disk. Read when the detector is loaded; never modified.
\end{itemize}

\subsubsection{Access Routine Semantics}

\noindent \texttt{extractImage}():
Require a supported image from M4's input contract. Return a Signals record containing visual observations and model inferences separately from camera metadata (Exchangeable Image File Format, EXIF) and Content Credentials (Coalition for Content Provenance and Authenticity, C2PA) checks. Use the separate provenance fields defined in Table~\ref{TblTypes}; do not infer deliberate stripping from absence. Treat missing metadata as a limited observation and never treat EXIF presence as verified authorship. If required parsing or model work fails, identify the unavailable dependency so the coordinator can apply its failure policy.

%% ------------------------------------------------------------------ M9
\clearpage
\section{MIS of M9 Assessment policy} \label{sec:m9}

M9 applies a frozen decision profile and abstains when a directional assessment is unsupported.

\subsection{Module}

Assessment policy

\subsection{Uses}

No design module.

\subsection{Syntax}

The following constants and operations form the public interface. Shared initialization is defined in Section~\ref{sec:notation}.

\subsubsection{Exported Constants}

None. Thresholds are validation-derived fields of the profile.

\subsubsection{Exported Access Programs}

\begin{center}
\begin{tabular}{p{3.3cm} p{4cm} p{4cm} p{2.4cm}}
\hline
\textbf{Name} & \textbf{In} & \textbf{Out} & \textbf{Exceptions} \\
\hline
\texttt{initialize} & $p:\ty{Config}$ & $\ty{Unit}\mid\ty{Reason}$ & None \\
\texttt{assess} & $q:\ty{Signals},\ p:\ty{Config}$ & $\ty{Decision}\mid\ty{Reason}$ & None \\
\hline
\end{tabular}
\end{center}

\subsection{Semantics}

The state and environment below are private; callers rely on the operation effects.

\subsubsection{State Variables}

$p_9:\ty{Config}$ is immutable after initialization and identifies compatible calibration/decision profiles by model and modality. $D_9:\operatorname{Optional}(\ty{ModelArtifact})$ holds loaded calibration parameters and decision tables, absent before initialization. No request state is retained.

\subsubsection{Environment Variables}

\begin{itemize}
\item $\mathit{profileFile}:\ty{FileSet}$: the released calibration and decision profile file on the service host's disk. Read when the profile is loaded; never modified.
\end{itemize}

\subsubsection{Access Routine Semantics}

\noindent \texttt{assess}():
Verify model version, applicable task, required signal availability and each evidence claim's support. If a binary score is available, let $\hat p$ denote its calibrated synthetic/manipulation score and $0<\tau\le0.5$ the frozen abstention margin. A directional check completes only when the profile applies, there is at least one supported observed signal, the signals agree and $|\hat p-0.5|\ge\tau$. The lower-signal direction additionally requires $\hat p\le0.5-\tau$. Return inconclusive when a condition is absent or cannot be evaluated. Opinion or satire alone has no bearing on truth. Use domain support and signal agreement to decide whether the generation task applies.

Let $b$ mean that the profile matches, applicability is supported, all required signals are available and non-conflicting, at least one observed signal is grounded, and the score is defined. Then
\[
out:=\begin{cases}
\text{completed, lower direction},& b\land\hat p\le0.5-\tau,\\
\text{completed, higher direction},& b\land\hat p\ge0.5+\tau,\\
\text{inconclusive},&\text{otherwise}.
\end{cases}
\]
Undefined scores are never compared numerically. Evidence text contains safe descriptors and offsets or external locators, never raw selected excerpts that would survive cleanup.

Every completed or inconclusive decision includes the coded limitation \texttt{LIMIT\_\allowbreak{}NO\_\allowbreak{}GUARANTEE}, rendered as ``This assessment does not guarantee authenticity or prove fraud,'' plus task-specific limitations and an independent next step. An abstention may display safe supporting observations. Missing provenance alone cannot drive a higher-signal decision. A dependency error reaches M5 separately from this decision. A score describes an assessment of the selected item within the detector's task.

%% ------------------------------------------------------------------ M10
\clearpage
\section{MIS of M10 Offline evaluation} \label{sec:m10}

M10 evaluates licensed benchmark inputs independently of live browsing.

\subsection{Module}

Offline evaluation

\subsection{Uses}

M4, M7--M9.

\subsection{Syntax}

The following constants and operations form the public interface. Shared initialization is defined in Section~\ref{sec:notation}.

\subsubsection{Exported Constants}

SRS evaluation targets \Const{MIN\_\allowbreak{}F1\_\allowbreak{}SCORE}, \Const{MIN\_\allowbreak{}COVERAGE}, \Const{MAX\_\allowbreak{}FPR},\newline \Const{MIN\_\allowbreak{}BASELINE\_\allowbreak{}DELTA}, \Const{MAX\_\allowbreak{}REPRO\_\allowbreak{}DELTA}.

\subsubsection{Exported Access Programs}

\begin{center}
\begin{tabular}{p{3.3cm} p{4cm} p{4cm} p{2.4cm}}
\hline
\textbf{Name} & \textbf{In} & \textbf{Out} & \textbf{Exceptions} \\
\hline
\texttt{initialize} & $p:\ty{Config}$ & $\ty{Unit}\mid\ty{Reason}$ & None \\
\texttt{evaluate} & $m:\ty{Manifest},\ p:\ty{Config}$ & $\ty{Run}\mid\ty{Reason}$ & None \\
\texttt{recompute} & $r:\ty{Run}$ & $\ty{MetricReport}\mid\ty{Reason}$ & None \\
\hline
\end{tabular}
\end{center}

\subsection{Semantics}

The state and environment below are private; callers rely on the operation effects.

\subsubsection{State Variables}

$r:\operatorname{Optional}(\ty{Run})$, initially absent; $p_{10}:\ty{Config}$, fixed at initialization. The run contains only licensed benchmark identities and outcomes, never live browse-session state.

\subsubsection{Environment Variables}

\begin{itemize}
\item $\mathit{datasetFiles}:\ty{FileSet}$: the versioned labelled benchmark files on the developer host's disk. Read by \texttt{evaluate}; never modified.
\item $\mathit{outputDir}:\ty{Directory}$: the experiment output directory on the developer host's disk. Changed by \texttt{evaluate}, which writes the run record, and by \texttt{recompute}, which writes the metric report.
\end{itemize}

\subsubsection{Access Routine Semantics}

\noindent \texttt{evaluate}():
Validate the licensed corpus revision, labels and input eligibility. Assign related sources and duplicate groups to one split. Freeze exact/near-duplicate checks, seed, split identities, subgroup tags, model versions and validation thresholds before held-out evaluation. Apply live modality/assessment contracts and record exactly one outcome per eligible item, including abstentions and errors. Set $r'$ to this run, write its artifact in the configured output directory, and return $out:=r'$. Never access live selections.

\medskip\noindent \texttt{recompute}():
For binary ground-truth label $y_i$ and directional prediction $\hat o_i$, let $I$ be all eligible held-out IDs, $A$ the IDs with completed directional assessments, $N_0=|\{i\in I:y_i=0\}|$, $A_0=|\{i\in A:y_i=0\}|$, and $\mathrm{FP}=|\{i\in A:y_i=0\land\hat o_i=1\}|$. Compute assessed-only macro-average $F_1$ and the baseline delta on the same assessed IDs. Compute coverage over all eligible items and false-positive rates with explicit denominators, following SRS NFR1: $\mathrm{Coverage}=|A|/|I|$, $\mathrm{FPR}_{\mathrm{assessed}}=\mathrm{FP}/A_0$ and $\mathrm{FPR}_{\mathrm{full}}=\mathrm{FP}/N_0$. Pass only if assessed macro-average $F_1\ge\Const{MIN\_\allowbreak{}F1\_\allowbreak{}SCORE}$, baseline improvement $\ge\Const{MIN\_\allowbreak{}BASELINE\_\allowbreak{}DELTA}$ on the same $A$, coverage $\ge\Const{MIN\_\allowbreak{}COVERAGE}$ and both FPRs $\le\Const{MAX\_\allowbreak{}FPR}$. An empty required two-class test population or undefined assessed-negative FPR cannot pass. Authentic-only fairness slices report FPR and coverage and mark two-class $F_1$ not applicable. Separately report full-population diagnostics, abstention/error counts, slice and per-class coverage, and the observed environment. Recompute only from input $r$, without changing its outcomes, and write and return $out:=\ty{MetricReport}$. Signal an invalid manifest or missing outcomes instead of reporting an acceptance pass.

For each fixed class $c\in\{0,1\}$, count $\mathrm{TP}_c$, $\mathrm{FP}_c$ and $\mathrm{FN}_c$ over $A$. Define
\[
 F_{1,c}=\begin{cases}
 \dfrac{2\mathrm{TP}_c}{2\mathrm{TP}_c+\mathrm{FP}_c+\mathrm{FN}_c},
 &2\mathrm{TP}_c+\mathrm{FP}_c+\mathrm{FN}_c>0,\\
 0,&\text{otherwise},
 \end{cases}
 \qquad F_1^{\mathrm{macro}}=\frac{F_{1,0}+F_{1,1}}{2}.
\]
Use the same fixed classes and convention for the baseline on identical assessed IDs; withholding a class cannot yield a perfect two-class score. Reject duplicate IDs, label/split mismatches and missing or extra outcomes before computing metrics.

\clearpage

\bibliographystyle {plainnat}
\bibliography {../../../refs/References}

\clearpage

\clearpage{}

\section{Appendix --- Generative AI Use Disclosure}

This section records the AI tools used to prepare and verify the Module Interface Specification.

\subsection{AI Tools Used}

\noindent OpenAI Codex (GPT-6.1), Anthropic Claude Sonnet 5.5, Google Gemini Flash 3.8, and Google Gemini 4 Hydrogen were used for drafting, type-checking, and verifying consistency across module interfaces and course checklists.

\subsection{How and Where Used}

\noindent Calvin Salsali used OpenAI Codex (GPT-6.1) to draft the shared types, M1--M3 browser interface signatures, and consent/display semantics.

\medskip\noindent Richard Han used Anthropic Claude Sonnet 5.5 to review M4 and M6 security, bounds, and ephemeral-cleanup contracts against the Hazard Analysis and MIS checklist.

\medskip\noindent Jaden Moore used OpenAI Codex (GPT-6.1) and Google Gemini Flash 3.8 to formalize common initialization rules, M5's 5-row lifecycle transition table, and typed state/environment definitions.

\medskip\noindent Raymond Tao used Google Gemini 4 Hydrogen to align M7--M10 signatures and recomputation math with the Fall ML Model Report and add the individual design reflections.

\subsection{Verification of AI-Generated Content}

\noindent Each module owner reviewed their interface signatures, state invariants, exported constants, and routine semantics line by line against the SRS, Hazard Analysis, V\&V Plan, and MIS checklist.

\medskip\noindent Automated specification checks passed via \texttt{python3 test/\allowbreak{}documentation/\allowbreak{}verify\_\allowbreak{}specification.py -{}-with-hazard} (8 tests) and \texttt{python3 test/\allowbreak{}documentation/\allowbreak{}verify\_\allowbreak{}design.py} (verifying all 10 modules, Uses edges, 16 exported constant literals, initialization rules, and all 32 Boolean guard assignments in Table~\ref{TblLifecycle}), and the compiled PDF was inspected for layout and references.

\newpage{}

\section{Appendix --- Reflection}

This appendix reflects on the Revision $-1$ module decomposition and interface specification across design prompts 1--6. Prompt 7 (implementing one module from Team 06's MIS) will be completed after the Revision 0 peer-implementation milestone in January 2027, linking that pull request directly here.

\subsection{Calvin Salsali}

What worked cleanly was keeping all browser-extension weirdness inside M1, M2, and M3 so the backend modules do not care about Manifest V3 service workers or tab lifecycles. The most useful design correction was adding \ty{ConsentPreview} to M2: our first pass promised a disclosure before consent without returning policy version or route eligibility in the signature, which would have forced M2 and M1 to guess each other's UI state. I pushed for keeping authoritative consent and result rendering in an extension-owned view instead of injecting an interactive overlay directly into arbitrary host webpages, because host page styles and scripts are untrusted and hard to fight cleanly in \textbf{SFWRENG 4HC3 (Human Computer Interfaces)} terms. The main limitation of this design is that we still cannot prove host-page overlays are completely spoof-proof if we ever mirror anything on-page, and we left UI component details out of this revision until we build the extension prototype.

\subsection{Raymond Tao}

Splitting signal extraction (M7/M8) from the decision policy (M9) and offline evaluation (M10) was the right call after seeing how badly our Fall ML pilot models behaved on Wikipedia and Tiny-GenImage. Because M7 and M8 just emit typed \ty{Signals} and separate \ty{Provenance} fields while M9 owns the abstention band and applicability check, swapping TF--IDF for RoBERTa later will not touch the extension or coordinator at all. The trickiest part to nail down was M10's \texttt{recompute} contract: drawing on \textbf{SFWRENG 4AL3 (Applied Machine Learning)}, \textbf{STATS 3Y03 (Probability and Statistics for Engineering)}, and our V\&V feedback, I wanted the exact selective macro-$F_1$, zero-denominator convention, and assessed-vs.-full FPR formulas written straight into the MIS so there is zero ambiguity when coding \texttt{src/benchmark/}. One limitation we noticed is that an absent genre/applicability gate forces M9 to fail closed to \texttt{inconclusive}, so getting a real held-out genre check working is now a clear blocker before Revision 0.

\subsection{Jaden Moore}

Writing the MIS before touching FastAPI endpoints forced us to solve concurrency races on paper first. Giving M5 single ownership of service request state while M3 owns the active tab view means a stale response arriving late---or right as the user hits cancel---gets dropped cleanly by key matching without locking the UI. Our biggest fix during review was adding explicit \texttt{initialize} guards across all stateful objects and separating the 1.0-second raw-content cleanup deadline from a 60-second safe status-delivery record age; otherwise an implementer would either leak raw payloads or forget late-status polling entirely. We explicitly rejected using Celery/Redis or a persistent database in Revision $-1$ (drawing on \textbf{SFWRENG 3BB4 (Concurrent System Design)} and \textbf{SFWRENG 3A04 (Software Design II)}): in-memory request records are far simpler to reason about and avoid creating accidental persistent stores of user text.

\subsection{Richard Han}

Turning our FMEA privacy controls into M6's \ty{WorkerLease} model (\texttt{hold}, \texttt{borrow}, \texttt{release}) made ephemeral retention concrete enough to actually test. At first M6 described a lease without any operation creating one, which an independent programmer could never have implemented consistently, so making \texttt{hold} issue the lease and \texttt{release} measure the 1.0-second deadline from the terminal decision timestamp fixed a real interface hole. The honest limitation of M6---which we made sure to state plainly rather than oversell---is that Python object-reference release is logical cleanup, not forensic heap wiping, and once a user approves sending text to an external API, M6 cannot un-send delivered bytes. Choosing strict request-scoped leases over global worker pools trades a bit of memory overhead for much cleaner \textbf{SFWRENG 3S03 (Software Testing)} verification when we test cancellation and exception cleanup.

\end{document}
